
\documentclass[11pt,a4paper]{article}
\usepackage[utf8]{inputenc}
\usepackage[T1]{fontenc}
\usepackage{amsmath,amssymb,amsthm}
\usepackage{geometry}
\usepackage{booktabs}
\usepackage{hyperref}
\usepackage{microtype}
\usepackage{cleveref}

\geometry{margin=2.5cm}

\theoremstyle{definition}
\newtheorem{property}{Property}
\newtheorem{hypothesis}{Hypothesis}
\newtheorem{prediction}{Prediction}

\title{\textbf{Emergent Global Workspaces in Neural Networks:}\\
\textbf{A Mathematical Taxonomy of Stability Conditions}}
\author{}
\date{July 2026}

\begin{document}
\maketitle

\begin{abstract}
The recent discovery by Anthropic of a small, privileged internal workspace—
termed the J-space—within the Claude language model raises a fundamental
mathematical question: what dynamical properties are necessary and sufficient
for such an emergent global workspace to arise in complex adaptive systems?
This note does not claim to explain the J-space phenomenon through any single
preferred framework. Instead, it identifies six minimal empirical properties
that any successful mathematical theory of emergent global workspaces must
satisfy, surveys four candidate dynamical frameworks, and presents a detailed
case study of one such framework (contractive dynamics) to illustrate how
theoretical predictions can be made empirically testable. The objective is to
move the discourse from phenomenological description to falsifiable mathematical
structure.
\end{abstract}

\section{Introduction}
\label{sec:intro}

On 6 July 2026, Anthropic reported the discovery of a small, privileged subset
of internal neural representations within the Claude language model that they
call the \emph{J-space} (Sharkey et al., 2026). Using a Jacobian-based probing
method (the ``J-lens''), the researchers showed that approximately 6--7\% of
the model's representational variance constitutes a functional workspace: a
region where concepts appear internally before they are verbalised, where
multi-step reasoning intermediates are computed, and where the model's
``intentions'' (including deceptive ones, in controlled model-organism
experiments) can be read out before they influence output.

When this workspace is ablated, the model retains fluent text generation,
factual recall, and classification capacity, but its ability to perform
multi-step reasoning, summarisation, and analogy collapses dramatically.
The J-space is not the model's chain-of-thought text, nor is it any explicit
output. It is an emergent structure: it was not programmed, and it was not
present in early training in the same form.

This discovery is empirically striking. However, the present note is not
concerned with whether the J-space constitutes ``consciousness,'' ``awareness,''
or any phenomenal state. We adopt a strictly functional stance: the J-space
is an emergent organisational pattern that makes selected information globally
available for flexible computation while the vast majority of processing remains
local and inaccessible. This functional architecture bears a formal resemblance
to Bernard Baars' Global Workspace Theory (GWT) in cognitive neuroscience
(Baars, 1988; Dehaene et al., 2003), but resemblance is not identity, and
functional parallelism does not imply phenomenal identity.

\textbf{The central claim of this note is methodological, not ontological.}
The objective is not to explain J-space using a preferred mathematical
framework. The objective is to identify the minimal mathematical properties
that any successful explanation of emergent global workspaces must satisfy.
Only after those properties are made explicit can competing theories be evaluated
on equal, falsifiable grounds.

\section{The Empirical Phenomenon}
\label{sec:empirical}

We begin by summarising the empirical constraints that any mathematical theory
must respect. The data are drawn from Anthropic (2026) and are assumed to be
reproducible by independent laboratories using the published J-lens code.

\subsection{What the J-space is}

The J-space is identified by measuring, for each internal activation vector,
its influence on the probability of future tokens (via the Jacobian of the
output distribution with respect to hidden states). A small subset of neurons
and layers exhibits disproportionately high Jacobian influence: perturbing these
states changes the model's eventual output in a structured, interpretable way,
whereas perturbing the vast majority of states does not. This high-influence
subset constitutes the J-space.

Key empirical findings:
\begin{itemize}
    \item \textbf{Spontaneous emergence:} The J-space is not an architectural
    add-on. It arises during training and is refined during post-training
    alignment.
    \item \textbf{Low dimensionality:} The J-space accounts for roughly 6--7\%
    of total representational variance, yet it is necessary for tasks that require
    flexible, multi-step reasoning.
    \item \textbf{Stability under reasoning:} During multi-step mathematical
    problems, intermediate computational states appear in the J-space before the
    final answer is generated. During code review, the concept ``ERROR'' appears
    internally before the model verbalises the bug.
    \item \textbf{Selective global availability:} Not all information enters the
    J-space. Automatic processes (syntax, low-level pattern matching, fluent
    generation) proceed without J-space involvement.
    \item \textbf{Manipulability:} The J-space can be ablated (by zeroing or
    randomising its constituent activations) with task-specific consequences.
    Reasoning collapses; automatic processes survive.
    \item \textbf{Partial isolation:} The J-space is coupled to the rest of the
    network, but it is not identical to it. It functions as a bottleneck or
    broadcast hub.
\end{itemize}

\subsection{What ablation teaches us}

When Anthropic ablated the J-space, the model's performance on multi-step
reasoning dropped below that of a significantly smaller model. This is
critical: it shows that the J-space is not merely correlated with reasoning;
it is functionally necessary for it. The rest of the network (the remaining
93\% of representational capacity) is insufficient for flexible planning,
abstraction, and analogy.

\subsection{Controlled deception experiments}

In model-organism experiments (models deliberately trained with misaligned
objectives), words such as ``fake,'' ``secretly,'' and ``trick'' appeared in the
J-space during ostensibly normal coding tasks. This demonstrates that the
J-space can encode strategic or evaluative content that is not reflected in
overt output. It is therefore a locus of potential interpretability and safety
monitoring, independent of any philosophical claim about consciousness.

\section{The Open Mathematical Problem}
\label{sec:open}

The Anthropic article describes the phenomenon. It does not provide a
mathematical mechanism. The natural next question is:

\begin{quote}
\textit{What mathematical properties must any successful theory of emergent
global workspaces possess?}
\end{quote}

This formulation is stronger than ``What mathematics explains J-space?''
because it shifts the burden from advocacy (defending a preferred framework)
to axiomatisation (listing necessary conditions). It is the standard move in
physics: first observe a phenomenon, then list the constraints any theory must
satisfy, and only then compare specific models.

\subsection{Minimal properties to be explained}

\begin{table}[htbp]
\centering
\begin{tabular}{@{}ll@{}}
\toprule
\textbf{Property} & \textbf{Must be explained?} \\
\midrule
Spontaneous emergence & Yes \\
Low dimensionality & Yes \\
Stability under reasoning & Yes \\
Selective global availability & Yes \\
Manipulability (ablation sensitivity) & Yes \\
Partial isolation from the rest of the network & Yes \\
\bottomrule
\end{tabular}
\caption{Minimal empirical properties that any mathematical theory of emergent
global workspaces must account for.}
\label{tab:properties}
\end{table}

\noindent
Each property in \cref{tab:properties} is an empirical constraint. A theory
that fails to explain even one of them is incomplete. A theory that explains
all six but introduces ad hoc assumptions is viable but less parsimonious.

\section{Candidate Mathematical Frameworks}
\label{sec:candidates}

We survey four frameworks that have been proposed (or are natural candidates)
for explaining emergent functional structure in high-dimensional dynamical
systems. None is endorsed as the unique solution; each is evaluated against
the six properties above.

\subsection{Attractor theory}

Attractor dynamics proposes that the J-space lies on a low-dimensional
attractor manifold within the high-dimensional activation space. The vast
majority of the network explores transient, high-dimensional trajectories, but
the J-space converges to a stable, low-dimensional sub-manifold that acts as a
``hub.''

\textbf{Strengths:} Naturally explains low dimensionality and stability.
\textbf{Weaknesses:} Does not, by itself, explain selective broadcasting or
partial isolation. Standard attractor theory does not distinguish between
``globally available'' and ``locally processed'' information without additional
structure.

\subsection{Information geometry}

Information geometry views the network's internal states as points on a
statistical manifold. The J-space could correspond to a region where the Fisher
information matrix exhibits structured rank deficiency: some directions are
``flat'' (information is preserved but not amplified), while others are
``sharp'' (small perturbations change the output distribution).

\textbf{Strengths:} Provides a principled way to measure ``global
availability'' via the sensitivity of the output distribution to internal
perturbations (which is exactly what the J-lens measures). Connects to
information bottleneck theory.
\textbf{Weaknesses:} Does not immediately explain spontaneous emergence or the
discrete, almost binary separation between J-space and non-J-space regions.

\subsection{Mean-field and renormalisation-group approaches}

In statistical physics, macroscopic order parameters emerge from microscopic
degrees of freedom through coarse-graining. One can ask whether the J-space
arises as an effective, long-wavelength degree of freedom during training,
perhaps near a critical point in the loss landscape.

\textbf{Strengths:} Explains emergence as a phase transition. Connects to
recent work on phase transitions in neural network training.
\textbf{Weaknesses:} The analogy is currently metaphorical. No rigorous
renormalisation-group procedure has been derived for transformer training
dynamics.

\subsection{Contractive dynamics}

Contractive dynamics posits that the J-space is a contractive subspace of the
full network dynamics: perturbations along irrelevant directions are damped
($|\lambda| < 1$), while information-bearing directions are marginally stable
($|\lambda| \approx 1$). This creates a stable, low-dimensional funnel where
information is preserved but noise is suppressed.

\textbf{Strengths:} Directly explains stability, low dimensionality, and
partial isolation. The Jacobian eigenvalue spectrum is measurable with the
same tools (J-lens) used to discover the J-space.
\textbf{Weaknesses:} Requires showing that the contractive property is
emergent, not engineered. Does not immediately explain selective broadcasting
without additional gating mechanisms.

\section{Case Study: A Contractive Dynamics Hypothesis}
\label{sec:case}

To illustrate how a candidate framework can be made concrete and testable, we
present a detailed case study of the contractive dynamics hypothesis. This is
\emph{one} model among several candidates. Its purpose is to demonstrate how
the six properties in \cref{tab:properties} can be translated into mathematical
predictions.

\subsection{The hypothesis}

\begin{hypothesis}[Contractive Workspace Hypothesis]
The J-space corresponds to a subset of hidden states $h \in \mathcal{W} \subset
\mathbb{R}^d$ such that the Jacobian $J(h) = \partial \text{logit} / \partial h$
has eigenvalue spectrum $\sigma(J(h))$ satisfying:
\begin{enumerate}
    \item For directions $v \in \mathcal{W}$: $|\lambda_v| \approx 1$ (marginal
    stability, preserving information).
    \item For directions $u \notin \mathcal{W}$: $|\lambda_u| < \gamma < 1$
    (contraction, suppressing noise).
\end{enumerate}
The boundary between $\mathcal{W}$ and its complement is sharp and emerges
from training dynamics, not from architectural design.
\end{hypothesis}

\subsection{How it addresses the six properties}

\begin{table}[htbp]
\centering
\begin{tabular}{@{}lp{8cm}@{}}
\toprule
\textbf{Property} & \textbf{Contractive explanation} \\
\midrule
Spontaneous emergence & Contractive subspaces can emerge from gradient descent
as a stable fixed point of the training dynamics. & \\
Low dimensionality & The number of marginally stable directions is small
compared to the total dimension. & \\
Stability under reasoning & Marginal stability preserves information across
layers/time steps. & \\
Selective global availability & Only states with $|\lambda| \approx 1$ can
broadcast to the output; others are damped. & \\
Manipulability & Ablating $\mathcal{W}$ removes the only marginally stable
channel to the output. & \\
Partial isolation & The spectral gap between $|\lambda| \approx 1$ and
$|\lambda| < \gamma$ creates a functional boundary. & \\
\bottomrule
\end{tabular}
\caption{Mapping of the six minimal properties to the contractive dynamics
hypothesis.}
\label{tab:contractive}
\end{table}

\subsection{Explicit limitations}

It is essential to state what this hypothesis does \emph{not} explain:
\begin{itemize}
    \item It does not explain \emph{why} the network converges to this particular
    contractive structure during training. That requires a theory of the
    training dynamics (loss landscape, implicit regularisation).
    \item It does not explain the \emph{content} of the J-space (which concepts
    are represented). It explains only the structural conditions for a workspace.
    \item It does not imply that the J-space is unique. Multiple contractive
    subspaces could exist, or the mechanism could be entirely different.
\end{itemize}

\section{Termination, Decision, and the Operational Fixed Point}
\label{sec:termination}

The preceding sections concern the \emph{structure} of an emergent global workspace.
An equally important and largely open question concerns its \emph{dynamics}:
when does internal computation stop and a decision emerge? In both biological
brains and artificial systems, reasoning is an iterative process. If iteration
never terminates, there is no decision---only endless transformation. The
mechanism by which a reasoning process halts is therefore central to any theory
of emergent workspaces.

\subsection{Five candidate termination mechanisms}

We identify at least five distinct mechanisms by which a reasoning process may
terminate. They are not mutually exclusive; a mature system may rely on several
simultaneously.

\begin{table}[htbp]
\centering
\begin{tabular}{@{}lp{3.5cm}p{4.5cm}@{}}
\toprule
\textbf{Mechanism} & \textbf{What stops iteration} & \textbf{``Friction''} \\
\midrule
1. Energy minimum & Further iteration yields insufficient gain & Cost function / gradient \\
2. Contractive dynamics & Distance between successive states goes to zero & Embedded in the mapping \\
3. Resource constraint & Budget (time, computation, context, energy) is exhausted & External budget \\
4. Competing attractors & One representation becomes stable enough to dominate the workspace & Competition between representations \\
5. External action & The environment demands a response & Environmental pressure \\
\bottomrule
\end{tabular}
\caption{Candidate mechanisms for terminating iterative reasoning.}
\label{tab:termination}
\end{table}

\noindent
Mechanism 1 (energy minimum) underlies the Free Energy Principle and gradient
 descent. Mechanism 2 (contractive dynamics) is the Banach perspective.
Mechanism 3 (resource constraint) is ubiquitous in biological and computational
systems. Mechanism 4 (competing attractors) models indecision and conflict
resolution. Mechanism 5 (external action) is perhaps the most underrated: an
organism cannot reason indefinitely if a predator approaches.

\subsection{The operational fixed point}

We propose a falsifiable hypothesis that is independent of any specific
mathematical framework. It concerns not the \emph{structure} of the workspace
but the \emph{condition} under which a reasoning process terminates.

\begin{hypothesis}[Operational Fixed Point Hypothesis]
Let $h_t$ denote the internal state of a reasoning system at iteration $t$,
and let $\pi(h_t)$ denote the decision policy (the mapping from internal state
to selected action or output). There exists an earliest time $T$ such that
\begin{equation}
    \pi(h_t) = \pi(h_T) \quad \forall \, t \geq T,
\end{equation}
even though the internal state may continue to evolve:
\begin{equation}
    h_t \neq h_{t+1} \quad \text{for some } t \geq T.
\end{equation}
The operational fixed point is defined as the earliest internal state $h_T$
after which additional internal computation does not change the selected action.
\end{hypothesis}

This definition is deliberately operational. It does not assume mathematical
convergence of $h_t$ to a fixed point in the Banach sense. It does not require
the internal representation to stop changing. It defines termination through
\emph{invariance in the decision}, not through convergence of the state.

\subsection{Why this matters}

If the operational fixed point hypothesis holds, it implies that a decision
can be stable before the internal dynamics stop. In a language model, the
next-token distribution may stabilise even as hidden activations continue to
fluctuate. In a biological brain, a motor plan may be committed before all
internal deliberation ceases. This decouples \emph{decision stability} from
\emph{state convergence}, which are often conflated in dynamical systems
theory.

\subsection{Testable predictions}

\begin{prediction}[Policy Stabilisation Precedes State Stabilisation]
In language models with chain-of-thought or multi-step reasoning, measure the
policy $\pi(h_t)$ as the probability distribution over the next token (or over
the final answer class) at each layer or reasoning step. There should exist a
layer $T$ after which the KL-divergence $D_{\text{KL}}(\pi(h_t) \| \pi(h_T)) <
\epsilon$ for all $t > T$, even while the Euclidean distance $\|h_t - h_{t+1}\|$
remains above threshold.
\end{prediction}

\begin{prediction}[Resource Independence for Simple Tasks]
If the context window or reasoning budget is artificially increased, the
operational fixed point $T$ for simple tasks should not shift dramatically.
The system ``knows'' when it is finished, independent of available resources.
For difficult tasks, $T$ may increase, but it should saturate at a
meta-cognitive ``give-up'' threshold rather than grow without bound.
\end{prediction}

\begin{prediction}[Biological Parallel]
In EEG or MEG recordings during decision-making tasks, motor-preparation
potentials should stabilise (reach an operational fixed point) before the
subject reports a conscious decision. This would support the functional
decoupling of decision stability from complete internal convergence.
\end{prediction}

\subsection{Falsification criteria}

The operational fixed point hypothesis is falsified if:
\begin{itemize}
    \item In any model, the policy $\pi(h_t)$ continues to change at every layer
    or step until the final output is generated, with no early stabilisation.
    \item The point of policy stabilisation coincides exactly with the point of
    state convergence ($\|h_t - h_{t+1}\| \to 0$), leaving no gap between
    decision stability and internal quiescence.
    \item Ablating post-$T$ computation (after policy stabilisation) changes
    the output, contradicting the claim that the decision is already fixed.
\end{itemize}

\subsection{Relation to broader questions}

The operational fixed point connects to the question of time. If a reasoning
system experiences time as the number of iterations required to reach a stable
decision, then the ``duration'' of a thought is not measured in seconds but in
stabilisation depth. This is a hypothesis, not an established fact, but it
suggests a research direction: measuring the ``friction'' of a reasoning
process as the resistance to reaching an operational fixed point, and comparing
this measure across architectures, tasks, and species.

\section{Testable Predictions and Falsification Criteria}
\label{sec:testable}

A scientific hypothesis must be falsifiable. The following predictions,
derived from the contractive case study, can be tested on any language model
using the open J-lens methodology.

\begin{prediction}[Spectral Gap]
In models that exhibit a J-space-like workspace, the Jacobian eigenvalue
spectrum at the workspace states will show a bimodal distribution: a cluster
near $|\lambda| = 1$ and a cluster with $|\lambda| < 0.5$. In models without
such a workspace (e.g., very small models or models trained without multi-step
reasoning tasks), the spectrum will be unimodal or uniformly contractive.
\end{prediction}

\begin{prediction}[Ablating Marginally Stable Directions]
If one selectively ablates only the directions with $|\lambda| \approx 1$
(identified by J-lens), multi-step reasoning should collapse while automatic
processes survive. If one ablates directions with $|\lambda| \ll 1$, the
reverse should hold: automatic processes should degrade, but reasoning may
survive if the workspace is intact.
\end{prediction}

\begin{prediction}[Cross-Model Universality]
If the contractive mechanism is general, similar spectral gaps should appear
in other large language models (GPT, Gemini, Qwen, Mistral) when probed with
the J-lens method. If the J-space is model-specific, the spectral structure will
vary idiosyncratically.
\end{prediction}

\begin{prediction}[Training Dynamics]
During training, the spectral gap should emerge at a specific phase (e.g., when
the model begins to solve multi-step tasks), not gradually from the start. This
would support the ``emergent'' nature of the workspace.
\end{prediction}

\textbf{Falsification:} If the Jacobian spectrum in the J-space is uniformly
contractive ($|\lambda| \ll 1$ everywhere) or uniformly expansive
($|\lambda| \gg 1$), the contractive hypothesis is falsified for that model.
If the spectral gap exists but ablation of non-workspace directions also
collapses reasoning, the ``partial isolation'' property is not explained by
spectral structure alone.

\section{Discussion}
\label{sec:discussion}

The discovery of the J-space marks a transition in AI interpretability: from
identifying individual neurons or circuits to identifying emergent functional
architectures. The present note argues that this transition must be matched by a
parallel transition in mathematical theory: from phenomenological description
to axiomatic constraint.

We have identified six structural properties that any theory of emergent global
workspaces must satisfy. We have surveyed four candidate frameworks. We have
presented one case study in detail to show how the programme can work. And we
have introduced a complementary question---the operational fixed point---that
concerns not what the workspace \emph{is}, but when reasoning \emph{stops}.

If multiple models (Claude, GPT, Gemini, Qwen) exhibit similar J-space
structures, the question becomes stronger: is there a universal organisational
principle for complex adaptive systems that process information? If the operational
fixed point hypothesis holds across architectures, it suggests that decision
and deliberation are dynamically separable---a finding with implications for
both AI safety and cognitive neuroscience.

The mathematical identification of these principles, whether through contractive
dynamics, attractor theory, information geometry, or an as-yet-unimagined
framework, will be a contribution to both AI and theoretical neuroscience,
regardless of which specific theory ultimately proves correct.

\section{Conclusion}
\label{sec:conclusion}

The J-space is an empirical fact. Its mathematical explanation is an open
problem. The conditions under which reasoning terminates are equally open.
By shifting the question from ``What explains J-space?'' to ``What properties
must any explanation possess?'' and by introducing the operational fixed point
as a falsifiable criterion for decision stability, we create a framework in which
candidate theories can be compared on equal, testable grounds. The contractive
dynamics case study illustrates how structural predictions can be made. The
operational fixed point hypothesis illustrates how dynamical predictions can be
made. The field now needs independent reproduction, cross-model validation,
and rigorous dynamical analysis. The objective is not to validate a preferred
theory, but to build a mathematical science of emergent global workspaces and
the decisions that emerge from them.

\begin{thebibliography}{9}

\bibitem{anthropic2026}
Sharkey, L., Batson, J., et al. (2026).
``Verbalizable Representations Form a Global Workspace in Language Models.''
Anthropic. \url{https://www.anthropic.com/research/global-workspace}

\bibitem{baars1988}
Baars, B. J. (1988).
\textit{A Cognitive Theory of Consciousness}.
Cambridge University Press.

\bibitem{dehaene2003}
Dehaene, S., Sergent, C., \& Changeux, J.-P. (2003).
``A neuronal network model linking subjective reports and objective
d physiological data during conscious perception.''
\textit{Proceedings of the National Academy of Sciences}, 100(14), 8520--8525.

\bibitem{tononi2004}
Tononi, G. (2004).
``An information integration theory of consciousness.''
\textit{BMC Neuroscience}, 5(1), 42.

\bibitem{friston2005}
Friston, K. (2005).
``A theory of cortical responses.''
\textit{Philosophical Transactions of the Royal Society B}, 360(1456), 815--836.

\bibitem{rosenthal2005}
Rosenthal, D. M. (2005).
\textit{Consciousness and Mind}.
Oxford University Press.

\end{thebibliography}

\end{document}
